\section{Reading the results}
\label{sec:output}

This section reads the output of use case~1 (\cref{sec:uc1}), produced by
\begin{shell}
mapping examples/uc1\_central\_sizing/application.json \textbackslash\\
\quad examples/uc1\_central\_sizing/configuration.json -{}-cores
\end{shell}

\subsection{What the tool understood}

The output starts with the application and the platform as the tool read
them: the number of task-chains and tasks, how many are bound to one ECU,
how many start with a sensor task and end with an actuator task, every ECU
with its cores, then the models and the method.

\listing[\fontsize{6.3pt}{7.6pt}\selectfont]{listings/uc1_input.txt}

\subsection{The summary table}
\label{sec:summary}

One column per allocation rule (here only FFDU, the rule of the file).

\begin{table}[h]
\centering\small\row
\caption{The rows of the summary table.}
\label{tab:summary-rows}
\begin{tabularx}{\textwidth}{@{}>{\raggedright\arraybackslash}p{50mm}X@{}}
\toprule\rowcolor{white}
row & meaning \\
\midrule
feasible & every sub-chain found a core; otherwise the line below the table says which one did not, and why (\cref{sec:infeasible}) \\
sub-chains, fragmentation rate, forced cuts & the partitioning: sub-chains, cuts in \% of the possible cuts, cuts between tasks bound to different ECUs \\
cores used $M$ & cores that run at least one task: \textbf{the sizing result}; then on central computers and on the other ECUs \\
lower bound on $M$ & no mapping of these sub-chains uses fewer cores (each ECU's bound load, and the total load, over the cap) \\
cores added, cores emptied & cores beyond the initial ones; cores emptied by the consolidation \\
schedulable cores SR & \% of cores where every task meets its deadline, after the priorities, the optional swaps, the relocations, and at the end \\
critical scaling factor & the margin of the platform (\cref{sec:concepts-metrics}); with sizing, the cap found and whether the target was reached \\
application / LET load, room left & in cores, summed over the cores used \\
task-chains schedulable / meeting deadline & \% of chains whose tasks are all schedulable; and also within their deadline over bounded routes \\
local search moves & moves kept by the deadline-driven local search \\
communications inside an ECU / inter-ECU & messages between two cores of one ECU (shared memory) and between ECUs (network flows) \\
replicated with FRER & flows sent over two routes \\
network delay bound, busiest port & the largest bound of a flow; the largest LET load of an output port \\
end-to-end response time, max reaction time & the largest values over the task-chains (\cref{sec:concepts-metrics}) \\
CPU time & time to compute the design point (and with the sizing) \\
\bottomrule
\end{tabularx}
\end{table}


\begin{minipage}[t]{0.47\textwidth}
\vspace{0pt}
\listingbox[\scriptsize]{listings/uc1_summary.txt}
\end{minipage}\hfill
\begin{minipage}[t]{0.5\textwidth}
\vspace{0pt}\small
\textbf{Partitioning.} 64 sub-chains for 73 tasks: most chains are cut
at both ends, where the sensor and the actuator are bound to zone
controllers (40 forced cuts).

\textbf{Cores.} 7 cores are used: 2 on the HPC and one in each zone
controller. The lower bound is 5: no mapping of these sub-chains can do
with fewer. No core was added; none had to be emptied.

\textbf{Schedulability.} Every core is schedulable after the priority
assignment; the relocation moves had nothing to do. The critical
scaling factor 1.302: every WCET may grow by 30\,\% before a core misses
a deadline.

\textbf{Load.} 2.70 cores of application load, 0.34 of LET copies, 2.56
cores left below the caps.

\textbf{Task-chains.} All on time; the local search kept 6 moves to get
there. 33 messages cross the network; the largest delay bound is
0.8\,ms on a port loaded at 0.5\,\%.

\textbf{Time.} The design point took less than a millisecond of CPU.
\end{minipage}


\subsection{Per ECU}

The second table gives, for every ECU and rule, the cores used, their
load in cores (LET included) and the smallest critical scaling factor of
its cores. Here \cmd{ZC\_FL} is the ECU with the least margin: its 1.302
is the factor of the platform.

\listing{listings/uc1_ecus.txt}

\subsection{Cores and task-chains}
\label{sec:cores-chains}

With \opt{cores}, a table per rule lists the cores in use: ECU, type and
speed, load and LET load (\% of the core), scaling factor,
schedulability, number of tasks, and the sub-chains on it as
\emph{chain.index}.

\listing[\tiny]{listings/uc1_cores.txt}

A second table follows every task-chain: its path as segments on cores,
with the indices of its sub-chains in brackets; the number of
communications inside an ECU and over the network; its maximum reaction
time and its deadline; and, if it is late, why: \emph{task late} (a task
misses its period), \emph{network} (a flow crosses an overloaded port) or
\emph{deadline} (the reaction time exceeds the deadline).

\listing[\fontsize{4.6pt}{5.6pt}\selectfont]{listings/uc1_chains.txt}

\begin{tip}[Reading a path]
\cmd{c6@ZC\_C(0) > c0@HPC(1) > c4@ZC\_RL(2)} (stability control) means:
sub-chain 0 runs on core c6 of zone controller \cmd{ZC\_C}, sub-chain 1
on core c0 of the HPC, sub-chain 2 on core c4 of \cmd{ZC\_RL}. The two
arrows are two inter-ECU communications, each costing LET delays and a
network transfer: 27\,ms of reaction time for tasks of 5\,ms.
\end{tip}

\subsection{When the mapping is infeasible}
\label{sec:infeasible}

A mapping is infeasible when a sub-chain fits on no core. The tool names
it and says why:
\begin{shell}
FFDU: infeasible, sub-chain 23.0 (0.8 \% of a reference core, 6144 KB,\\
free): not enough memory, at most 5120 KB left on the ECUs it may run on
\end{shell}
The other reasons are: no core of the ECUs it may run on has a type it
allows; the hardware is fixed (\opt{fixed-hardware}); the ECUs with memory
left already have their largest number of cores; or the sub-chain is
larger than an empty core (lower \opt{sigma}, or use faster cores).

\subsection{The result files}
\label{sec:outfiles}

\opt{out DIR} writes six files, one row per rule and item, so that a
spreadsheet or a script can take over.

\begin{table}[ht]
\centering\small\row
\caption{The files of \opt{out}.}
\label{tab:outfiles}
\begin{tabularx}{\textwidth}{@{}>{\raggedright\arraybackslash}p{24mm}X@{}}
\toprule\rowcolor{white}
file & content \\
\midrule
\cmd{summary.json} & one object per rule: the rows of the summary table, as numbers \\
\cmd{ecus.csv} & per rule and ECU: role, cores used, application and LET load (cores), critical scaling factor, memory used and capacity (KB), load of its egress and ingress network ports (\%) \\
\cmd{cores.csv} & per rule and core in use: ECU, type, speed, added or initial, load and LET load (\%), schedulable, critical scaling factor, tasks \\
\cmd{chains.csv} & per rule and task-chain: tasks, utilization, sub-chains, communications over the network and inside an ECU, end-to-end response time, maximum reaction time (both empty when the chain has no bound: a task above its period or an unbounded port), deadline, schedulable, network bounded, on time, why late \\
\cmd{mapping.csv} & per rule and task: chain and task (numbers and names), sub-chain, core, ECU, core type, priority, period, WCET on its core, response time (empty for a task above its period, which has none) \\
\cmd{system.json} & the timing model of each mapping (cores with their LET task and tasks, chains with the core of each task and the bound of each hop), for the full-system simulation (\cref{sec:ex-syssim}) \\
\bottomrule
\end{tabularx}
\end{table}
